\documentclass{article}
\usepackage[T1]{fontenc}
\usepackage{lmodern}
\usepackage{graphicx}
\usepackage{amsmath,amssymb}
\usepackage[a4paper,margin=2cm]{geometry}
\usepackage[absolute,overlay]{textpos}
\setlength{\TPHorizModule}{1mm}
\setlength{\TPVertModule}{1mm}
\usepackage[indonesian]{babel}
\usepackage{hyperref}
\setlength{\emergencystretch}{2em}
% unit: Halaman 7

\begin{document}

% === Halaman 7 (llm) ===

\begin{itemize}
    \item Perhatikan bahwa $a = 2$ bukan akar primitif dari 7 karena
    
    \[
    \begin{aligned}
    2^1 \pmod{7} &= 2; & 2^2 \pmod{7} &= 4; & 2^3 \pmod{7} &= 1 \\
    2^4 \pmod{7} &= 2; & 2^5 \pmod{7} &= 4; & 2^6 \pmod{7} &= 1
    \end{aligned}
    \]
    
    \item Nilai-nilai yang dihasilkan dari $2^1, 2^2, \dots, 2^6$ tidak semuanya berbeda dan tidak mencakup semua nilai di dalam modulus 7.
    
    \item Untuk menemukan semua akar primitif dari $p$, kita harus mencoba semua bilangan bulat dari $2, 3, \dots$
\end{itemize}

\end{document}
